\chapter{Case Studies}
\label{ch:cases}

The preceding chapters introduced tools for one question at a time. This chapter applies several of them together to a small number of well-known films, and to one synthetic scene. The claims made about each film are claims about what the model can state, given a feature of the film that is widely discussed. They are readings, offered as demonstrations of method, and none depends on attributing intentions to the makers.

\section{Rope: seamless composition with hidden seams}

\emph{Rope} (1948) was designed to appear as a continuous take, and is assembled from a small number of long takes joined by concealed cuts, many of them masked by a dark surface filling the frame. In the vocabulary of Chapter~\ref{ch:scenario}, the film is a word in the free monoid of shots whose seams have $\delta_\Img$ small while $\delta_\Sc$ need not be zero: the camera comes to rest against a dark back, the next take begins from the same configuration, and the two world states are close in appearance. The audience reads the sequence as a morphism in $\Path(\Sc)$, a single path, although it is a composite of path segments joined at matching seams.

The model separates two things that viewers may conflate. The film's effect of unbroken real time is the effect of $\delta_\Img$ being small at every seam. Its actual construction as a word of several shots is the fact that $\delta_\Sc$ need not vanish. By \Cref{prop:seam-lipschitz}, if the rendering is Lipschitz then small $\delta_\Sc$ suffices for small $\delta_\Img$, but the converse is the technique of the concealed cut: a large jump in world state hidden behind a rendering that cannot show it.

\section{Last Year at Marienbad: local consistency, global failure}

\emph{Last Year at Marienbad} (1961) is often described as one whose scenes, individually coherent, resist assembly into a single chronology. One reading in the present vocabulary treats each recalled scene as a local section of a sheaf of scene states over a cover of the film's story world, with each pair of overlapping recollections agreeing on what they share. If the family is compatible on overlaps and admits no global section, it is a contextual family in the sense of Chapter~\ref{ch:descent}. The reading makes a precise claim that can be argued against: that every pair, and indeed every small collection of scenes, admits a consistent state, while no state accounts for all. A critic who finds a consistent global chronology rejects the reading by exhibiting the global section. A critic who finds a pair of scenes that cannot both be true of one world rejects the claim of local consistency.

\section{Groundhog Day: appearance without history}

\emph{Groundhog Day} (1993) presents a day that repeats while one character retains knowledge across repetitions. As developed in \Cref{ex:groundhog}, the appearance is periodic while the state is not. The world clock obstruction of \Cref{prop:loop} applies to any attempt to give the film a global single-valued time with the image sequence as its only record. The obstruction disappears when the record of accumulated knowledge is placed in the history and the fold, which makes the world state non-periodic. This is the formal sense in which the film is about a history that the appearance does not show: two histories at the dawn of the first and tenth day have identical images and different admissible futures, which is \Cref{prop:appearance}.

The same example illustrates the preservation contracts of Chapter~\ref{ch:contracts}. A transformation that compresses the film by removing apparently repeated days would preserve event order only if every repetition were removable without effect, and fails to preserve the semantic fold, since the knowledge the protagonist accumulates is exactly the content of the repetitions.

\section{Rashomon: parallel accounts as distinct paths}

\emph{Rashomon} (1950) presents several accounts of one event. In the type-theoretic reading of \Cref{ex:rashomon}, the accounts are paths between the same two states of the world, and the film's concern is that they are not equal as paths even though they share endpoints. The reading separates the question of outcome, which is shared, from the question of how the outcome arose, which is contested. What a viewer must hold in mind is a family of paths and the claim, which the film neither confirms nor refutes, that no one of them is privileged.

\section{Perfect Days: a flat narrative potential}

\emph{Perfect Days} (2023) follows the repeated daily routine of a man in Tokyo, in a form that is repetitive without being a time loop. In the geometric vocabulary of Chapter~\ref{ch:scenario}, the narrative cost $U$ is nearly constant along the film's path, so the action functional is dominated by the perceptual term and critical paths are close to geodesics of the perceptual metric. The film's repetition is not a periodicity of the world clock but a recurrence of appearance with small, accumulating changes in state. Its interest, in this reading, lies in the history that the repeated appearances conceal, and in the small events that modify the fold without altering the image.

\section{A synthetic scene: one participant knows what the other does not}

A teaching scene that can be varied at will illustrates the machinery in one place. Two characters, $A$ and $B$, converse in a room. $A$ knows that the letter on the table is forged, and $B$ does not. The scene can be realized in several versions that differ in timing, soundtrack, cutting, and gesture.

\emph{Types.} The story state contains $\mathsf{Knows}(A,\mathsf{forged})$ and lacks $\mathsf{Knows}(B,\mathsf{forged})$. The event $\mathsf{reveal}(A,B)$ has the first as a precondition and, as an effect, adds the second. A version in which $B$ acts on the forgery before the reveal is ill typed.

\emph{State and appearance.} A shot of $A$ glancing at the letter has the same image whether or not $A$ knows. The state layer distinguishes the two cases, so by \Cref{prop:appearance} the image cannot determine the admissible futures, and the audience's inferences depend on what the history has established.

\emph{Transformation.} A compression that removes a pause before the reveal preserves event order. If the pause was an opportunity for the audience to infer $A$'s knowledge, it may fail a recovery or anticipation relation. A contract that lists only order permits the compression, and one that lists anticipation does not.

\emph{Equivalence.} Moving the reveal by $\Delta$ seconds changes comprehension for $\Delta$ above some threshold that depends on the viewer and task. By \Cref{prop:nontransitive}, a sequence of small shifts of the reveal can accumulate to a visible change although each step is below threshold.

\emph{Direction.} The two commands ``$A$ learns the letter is forged'' and ``$A$ confronts $B$'' do not commute (\Cref{ex:noncomm}). A director issuing them in the wrong order obtains an inadmissible version, and the system must either refuse or replan.

The scene is synthetic and its soundtrack and cutting can be varied without referring to commercial footage, which makes it suitable for experiments. Findings from such scenes are findings about the constructed variations, and they do not generalize automatically to existing films.

\section{The synthetic scene, analyzed by time}\label{sec:timeindexed}

The previous section applied each tool once. This section revisits the forged-letter scene segment by segment, using the specification of Section~\ref{sec:scene}, and for each segment records five things: what is \emph{observed} on screen and sound, how it is \emph{annotated} in the layers of Chapter~\ref{ch:layers}, the \emph{formal} object from earlier chapters that represents it, an \emph{alternative} edit and the relation it tests, and an \emph{evaluation} question that an experiment or reader could answer. Everything in the observation lines is stipulated by the specification, so the analysis is checkable against it. The evaluation lines are open.

\subsection*{$0$ to $6$ s: entrance}
\begin{description}
\item[Observed.] $A$ enters the kitchen. Ambient sound at $40$ dB.
\item[Annotated.] History: $\mathsf{enter}(A)$. State: $A$ knows the letter is forged. Appearance: nothing marks this. Audience: no knowledge of the forgery yet.
\item[Formal.] A seamless segment, a morphism in $\Path(\Sc)$; no seam. The state and appearance layers already differ here (\Cref{prop:appearance}).
\item[Alternative.] Cut the entrance in two shots. This introduces a seam and tests continuity in the sense of Chapter~\ref{ch:categories}.
\item[Evaluation.] Does a viewer perceive the second version as having a cut at the seam, given a match on $A$'s position?
\end{description}

\subsection*{$6$ to $18$ s: greeting}
\begin{description}
\item[Observed.] $B$ greets $A$ at $t=6$. Ordinary dialogue.
\item[Annotated.] History: $\mathsf{greet}(B,A)$. Neither character acts on the letter.
\item[Formal.] A word in the free monoid of shots. Cut density is a statistic of the word, not of the story (Chapter~\ref{ch:editing}).
\item[Alternative.] Replace three cuts by one: the cut rate drops, the events stay.
\item[Evaluation.] Is the perceived pace changed (the tolerance $\varepsilon$ of Chapter~\ref{ch:equivalence}) when the cut count changes and the events do not?
\end{description}

\subsection*{$18$ to $31$ s: the glance and the street}
\begin{description}
\item[Observed.] At $t=18$, $A$ glances at the letter. At $t=20$ to $22$ street noise at $64$ dB.
\item[Annotated.] The glance is an appearance with no history event; the street noise is a non-cue exceeding event.
\item[Formal.] The glance is the observable that supports $\mathsf{T}_{\mathrm{con}}$ over $\mathsf{T}_{\mathrm{int}}$ (Section~\ref{sec:tropeexample}); the street noise is a block of three exceeding steps (Section~\ref{sec:whisper}).
\item[Alternative.] Remove the glance (state unchanged, appearance changed), or lower the street noise by $1$ dB (occupation changed, no other relation).
\item[Evaluation.] Does removing the glance change which reading viewers report at $t=46$?
\end{description}

\subsection*{$31$ to $38$ s: the pause}
\begin{description}
\item[Observed.] A $7$ s pause with no dialogue and ambient sound only.
\item[Annotated.] Quiet interval. No history event.
\item[Formal.] The first moment where $\rho_{\mathrm{ant}}$ of Section~\ref{sec:altversions} bites: a quiet run of $7$ s against a threshold $\tau_0=5$ s. Parameter $p_1$ of Section~\ref{sec:perturb}.
\item[Alternative.] Shorten to $3$ s ($V_1$): order preserved, anticipation relation broken.
\item[Evaluation.] Between $3$ s and $7$ s, where does a viewer population's judgment of anticipation drop by more than $\varepsilon$, and is the dependence asymmetric as the model predicts?
\end{description}

\subsection*{$38$ to $58$ s: question, deflection, buzz}
\begin{description}
\item[Observed.] $B$ asks the question at $t=38$. $A$ deflects at $t=46$. $B$'s phone buzzes at $t=52$ at $54$ dB.
\item[Annotated.] History: $\mathsf{ask}$, $\mathsf{deflect}$, and the buzz event. The buzz is a causal cue.
\item[Formal.] The word $\mathsf{ask}\,\mathsf{deflect}$ is a prefix of a development of $\mathsf{T}_{\mathrm{con}}$, with the timing window checks of Section~\ref{sec:tropeexample}. The buzz fixes the cue limit $a_{\max}=4$ dB (\Cref{prop:cuesafe}).
\item[Alternative.] Reorder: place the audience disclosure at $t=18$ ($V_4$). The type of the question changes from mystery to irony.
\item[Evaluation.] Which of the two readings, $\mathsf{T}_{\mathrm{con}}$ or $\mathsf{T}_{\mathrm{int}}$, do viewers report at $t=46$, with and without the discriminating observation?
\end{description}

\subsection*{$58$ to $66$ s: the reveal and the reaction}
\begin{description}
\item[Observed.] $A$ reveals the forgery at $t=58$. A reaction shot of $B$ from $t=62$ to $66$.
\item[Annotated.] History: $\mathsf{reveal}$, then $\mathsf{react}$. The reveal adds $\mathsf{Knows}(B,\varphi)$. Audience: now informed.
\item[Formal.] The only step at which both the character-knowledge and audience-knowledge processes change. Montage surplus: the reaction shot's reading depends on the preceding line (Chapter~\ref{ch:editing}). Reaction duration is parameter $p_4$.
\item[Alternative.] Place the reaction before the reveal line, or lengthen it from $4$ s to $6$ s.
\item[Evaluation.] Is the reaction read as a response to the reveal in the reordered version? By how much does perceived salience change with duration?
\end{description}

\subsection*{$66$ to $78$ s: the aftermath}
\begin{description}
\item[Observed.] Dialogue and ambient only; a quiet run of $12$ steps.
\item[Annotated.] No history event occurs here; the segment carries the aftermath as appearance only.
\item[Formal.] Part of the $55$-step recovery interval between the street noise and the slam (Section~\ref{sec:whisper}), which is the longest quiet run in the scene and is untouched by any attenuation policy considered there.
\item[Alternative.] Insert an additional exchange ($+6$ s): changes running time and the placement of the slam.
\item[Evaluation.] Does the additional exchange change the perceived impact of the slam?
\end{description}

\subsection*{$78$ to $80$ s: the exit}
\begin{description}
\item[Observed.] $B$ leaves. Door slam at $72$ dB for $2$ steps.
\item[Annotated.] History: $\mathsf{exit}(B)$, a causal cue. State: $B$ is no longer present, so any further command that requires $B$ is inadmissible (Chapter~\ref{ch:linear}).
\item[Formal.] An exceeding block of two steps, and $\rho_{\mathrm{cue}}$, with $\rho_{\mathrm{sal}}(\kappa)$ for $\kappa>\mu$ in conflict with zero occupation (Section~\ref{sec:whisper}); a command \texttt{s d x} with its precondition (Section~\ref{sec:session}).
\item[Alternative.] $S^*$ (slam lowered by $9$ dB), or a soft slam of $51$ dB.
\item[Evaluation.] Do listeners prefer $S$ to $S^*$? Is the door slam still read as the end of the exchange at $63$ dB, and what salience margin $\kappa$ does the author intend?
\end{description}

\subsection*{$80$ to $96$ s: the music}
\begin{description}
\item[Observed.] Music enters at $t=80$ at $58$ dB and rises by $0.8$ dB per second to $70$ dB at $t=95$. The scene ends at $96$.
\item[Annotated.] Appearance only; no history event. A trailing exceeding block $t=87$ to $95$.
\item[Formal.] The largest block of occupation (nine of the $14$ exceeding steps). Music parameters $p_2,p_3$.
\item[Alternative.] Lower the music by $7$ dB (policy $S$) so that the block disappears.
\item[Evaluation.] What is lost: the modeled fit measure $q_2$ falls from $0.68$ to $0.27$. Is this reproduced by listeners?
\end{description}

The analysis shows how the tools divide the work. Only two events change the semantic state, the reveal and the exit. Pace statistics and sound levels describe appearance, tropes constrain state and audience knowledge together, and contracts name which of these a transformation may touch. The same segment appears in the analyses of several chapters, which is the purpose of carrying one scene through the book.

\section{A protocol for analyzing a real passage}\label{sec:protocol}

The film subsections of this chapter are readings, and none is a time-indexed analysis of footage. A reader who wants to carry out such an analysis on an actual film needs source material that this book does not supply, and should supply it with its provenance. The following protocol is the form that the synthetic analysis above takes when applied to a real passage. Section~\ref{sec:real} carries it out on a passage for which measurements exist.
\begin{enumerate}
\item \textbf{Source.} State the title, the edition or release, the runtime, and the player or medium. Timecodes belong to an edition and differ between releases, so they must be tied to one.
\item \textbf{Passage.} Give start and end timecodes, in a stated precision, for each segment analyzed.
\item \textbf{Observation.} For each segment, list what is seen and heard, in terms that another viewer can check without accepting the reading: the shots, the cuts, the dialogue, and the sound events, with timecodes.
\item \textbf{Annotation.} Assign each observation to a layer: history event, state fact, appearance, or audience knowledge. Mark which assignments are inferences.
\item \textbf{Claim type.} Mark each claim as an observation, a reading, or a hypothesis about perception, so that the evidence required differs correctly.
\item \textbf{Formal representation.} Name the object used (a path, a section of a sheaf, a schema instance, a command), and say which checks it passes. A representation that cannot fail a check is decoration.
\item \textbf{Alternative edit.} State a specific change, the relation it is meant to preserve, and the relation it breaks.
\item \textbf{Evaluation.} State the question a viewer study would answer, the population, and the tolerance.
\end{enumerate}
The protocol separates what can be verified from the film (items 1 to 3) from what is interpretive (4 to 6) and from what is empirical (7 and 8). A careful analysis keeps the three apart.

\section{A measured passage: \emph{Drive Through Fire}}\label{sec:real}

The synthetic scene is clear because it was built to be. To test whether that clarity made the framework look easier to apply than it is, the protocol of Section~\ref{sec:protocol} is applied here to a real passage for which measurements exist: the soundtrack analysis of the 2026 feature \emph{Drive Through Fire} in the author's monograph on acoustic occupation~\cite{flyxion2026occupation}. The numbers below come from two sources: Tables~7.3 to~7.5 and Section~7.5.7 of that monograph, and, for quantities the monograph does not report, a recomputation from its per-second table (\texttt{audio-analysis/results/seconds.csv} in the author's repository, commit \texttt{1c016ab}). The recomputation reproduces the monograph's nine occupation blocks exactly (the script \texttt{real\_passage.py} distributed with the source does this, and \texttt{verify.py} checks the transcribed tables for arithmetic consistency). The film itself, its picture, and its story were not available to this book. The recomputation shows that the tables follow from the per-second data, not that the per-second data follow from the film.

\paragraph{1. Source.} The 2026 release analysed in the monograph, runtime $1{:}27{:}02$. The level measure is ITU-R BS.1770 momentary loudness of the 5.1 mix, energy-averaged per second, in LUFS. The dialogue reference is $D=-34.6$ LUFS, the median level of the $1499$ seconds covered by subtitles ($799$ cues, applied offset $0.0$ s). Timecodes belong to this release.

\paragraph{2. Passage.} $1{:}00{:}39$ to $1{:}20{:}51$, a span of $1212$ s, bounded by the longest intervals without a qualifying recovery window.

\paragraph{3. Observation.} The monograph defines exceedance $E(t)=L(t)-D$, occupation as $E\ge\delta=10$ dB, and a recovery window as at least $5$ s at or below $D$. Terminology matters here: a \emph{recovery window} is a stretch of at least $5$ s at or below $D$, and what the monograph calls a recovery gap is the complementary interval \emph{without} a qualifying window, not a quiet interval. In the passage there are three qualifying windows, of $9$, $7$, and $13$ s, the last two separated by a single second above $D$.
\begin{center}
\small
\begin{tabular}{llr}
\toprule
interval & kind & seconds\\
\midrule
$1{:}00{:}39$ to $1{:}03{:}18$ & no qualifying window & $159$\\
$1{:}03{:}18$ to $1{:}03{:}27$ & recovery window & $9$\\
$1{:}03{:}27$ to $1{:}09{:}00$ & no qualifying window & $333$\\
$1{:}09{:}00$ to $1{:}09{:}07$ & recovery window & $7$\\
$1{:}09{:}07$ to $1{:}09{:}08$ & above $D$ & $1$\\
$1{:}09{:}08$ to $1{:}09{:}21$ & recovery window & $13$\\
$1{:}09{:}21$ to $1{:}20{:}51$ & no qualifying window & $690$\\
\bottomrule
\end{tabular}
\end{center}
The intervals sum to $1212$ s. The three windows total $29$ s, not the $30$ s stated in the monograph, whose figure is the length of the two short intervals between the three long ones ($9+21$ s) and so includes the $1$ s above $D$ and counts two windows where there are three. Counting every second at or below $D$, not only those in windows of at least $5$ s, there are $80$ such seconds in the passage, in $32$ separate runs, the longest $13$ s. The passage is therefore not free of quiet moments; its quiet moments are fragmented. These are acoustic opportunities for recovery under a stated rule, not measurements of a listener's recovery. Six occupation blocks (level-defined, merge gap $g=5$ s, minimum length $60$ s) lie inside the passage.
\begin{center}
\small
\begin{tabular}{llrrrrr}
\toprule
start & end & seconds & median $E$ & $Q_{95}(E)$ & throttle $T_b$ & subtitled s\\
\midrule
$1{:}00{:}48$ & $1{:}02{:}29$ & $101$ & $+16.2$ & $+19.1$ & $9.1$ & $10$\\
$1{:}04{:}24$ & $1{:}05{:}56$ & $92$ & $+16.8$ & $+20.7$ & $10.7$ & $0$\\
$1{:}06{:}46$ & $1{:}08{:}05$ & $79$ & $+16.9$ & $+22.8$ & $12.8$ & $10$\\
$1{:}09{:}24$ & $1{:}11{:}34$ & $130$ & $+12.5$ & $+16.5$ & $6.5$ & $48$\\
$1{:}13{:}45$ & $1{:}14{:}58$ & $73$ & $+15.5$ & $+18.8$ & $8.8$ & $13$\\
$1{:}15{:}09$ & $1{:}19{:}53$ & $284$ & $+16.0$ & $+20.4$ & $10.4$ & $30$\\
\bottomrule
\end{tabular}
\end{center}
Levels are in dB relative to $D$ and the throttle depth is $T_b=\max(0,Q_{95}(E\mid b)-\delta)$. The six blocks total $759$ s, which is $62.6\%$ of the $1212$ s of the span, the figure the monograph reports; before $1{:}00{:}39$ the share is $6.8\%$. Over the whole film, $30.8\%$ of audible seconds lie at least $10$ dB above $D$.

\paragraph{4. Annotation.} Everything above lies in the appearance layer of Chapter~\ref{ch:layers}, and specifically in sound. The monograph carries no shot boundaries and no history or state annotation for this passage: nothing says what happens in the story during any block. Subtitled seconds mark when speech occurs and not how loud it is, and the per-second level is the level of the whole mix, so a subtitled second inside a block records speech occurring in the loud material, not loud speech. The annotation step of the protocol therefore stops at the first layer, and every statement below about plot function would be an addition not supported by the source.

\paragraph{5. Claim types.} The table entries are observations (measured, in the monograph's grading). The monograph's statement that the loud plateau reflects compression or limiting is a hypothesis about production, graded as such there. The listener behaviour that motivated the study, turning the volume down for one to two minutes or more, is a single retrospective report without timestamps. It is the motivation of the study and not evidence in it.

\paragraph{6. Formal representation.} Three constructions of the earlier chapters apply, with differences that matter.

\emph{Exceedance and arrangement.} \Cref{prop:arrangement} says the histogram of levels does not determine the arrangement of exceedance. The real data show it concretely: the film-level figure of $30.8\%$ is the same whether occupation is spread evenly or concentrated, and here it is concentrated, with $13\%$ to $40\%$ of each ten-minute segment before $1{:}00{:}00$ and $62\%$ and $71\%$ in the next two. A further point is that the count of blocks depends on a parameter the raw definition of Section~\ref{sec:occupation} lacks. The monograph merges runs separated by at most $g=5$ s. The last two blocks are separated by $11$ s ($1{:}14{:}58$ to $1{:}15{:}09$). Rerunning the merge on the per-second occupied set gives six blocks inside the passage at $g=5$, five at $g=10$, and four at $g=11$ and $g=12$. The count does not fall only when the last two blocks merge, because merging also absorbs the short occupied episodes between blocks, which are not in Table~7.3: at $g=10$ the second and third blocks join across such episodes into one stretch from $1{:}03{:}44$ to $1{:}08{:}37$ (seconds $3824$ to $4117$), and at $g=11$ the last two blocks merge into one of $368$ s. The number of occupation blocks is therefore a statement about the exceedance set and a choice of $g$, and a report should give both.

\emph{Uniform attenuation.} The throttle depth $T_b$ is the uniform attenuation $a=T_b$ that brings the $95$th percentile of exceedance in block $b$ to the tolerance. By definition of the percentile, at least $95\%$ of the block's seconds then satisfy $E-a\le\delta$; the small residual fraction ($5.1\%$ to $6.5\%$ in the six blocks, recomputed from the per-second table) is built into the rule and is not an independent result. By \Cref{prop:mufflerizer} the exceedance sets are nested in $a$, so larger attenuations clear more. All such occupation figures use the \emph{fixed original} reference $D$. A uniform gain lowers dialogue and effects together, so measured against dialogue that has itself been lowered by $a$ the exceedance is unchanged, which is the listener's dilemma in the monograph. A reduction in occupation under a fixed reference is a statement about absolute level, and it says nothing about intelligibility or listening quality. The six values, $6.5$ to $12.8$ dB, are far above the $4$ dB cue-safe limit of the synthetic scene. That limit came from a stipulated floor on named cue sounds, and nothing in the source identifies cue sounds, so \Cref{prop:cuesafe} cannot be evaluated here. The quantity it needs, which sounds are the only evidence of a plot event, is the one the data lack.

\emph{Dialogue as a stem.} In the synthetic scene dialogue was a separate stem and was left unchanged. In the measurement it is not: the level is that of the whole mix. A global attenuation by $a$ lowers dialogue by $a$ as well, so $\rho_{\mathrm{dia}}$ fails in every block by between $6.5$ and $12.8$ dB. The stem-wise policy of Section~\ref{sec:whisper} needs a proxy stem. The monograph supplies one: subtitled seconds have a median centre-channel advantage of $16.0$ dB. The measured centre advantage shows how far this works, with a caution about what is measured. The seconds are subtitle-associated seconds, each a mixture of whatever sounds occur in that second, not isolated dialogue; channel location is not sound function. For subtitle-associated seconds over the whole film the median is $16.0$ dB, but \emph{inside the blocks} it is $+8.8$ dB at the median (quartiles $+3.9$ and $+12.0$, tenth percentile $-0.5$), against $+17.0$ dB outside them. Dialogue inside loud material is much less centre-dominant than dialogue elsewhere. Under the idealization that mix power is the sum of centre power $C$ and the rest $N$ with advantage $x=C-N$ in dB, removing $N$ entirely changes the mix level by
\[
10\log_{10}\bigl(1+10^{-x/10}\bigr)\ \mathrm{dB}\quad\text{(a decrease)},
\]
which is $0.11$ dB at $x=16$ but $0.54$ dB at $x=8.8$ and $1.5$ dB at $x=3.9$. Applied second by second to the $128$ subtitled seconds inside blocks ($8.5\%$ of all subtitled dialogue), the median decrease is $0.54$ dB and one in ten of these seconds falls by $3.3$ dB or more. The mix level of those seconds therefore contains a substantial non-centre part. Removing the non-centre channels might remove masking material, useful speech, or both, and the levels cannot say which. The reported changes therefore describe the channel mixture under the stated idealization, and dialogue preservation is left unresolved. (The idealization is approximate: the pipeline's RMS level of the mix is, on median, $3.0$ dB below the power sum of its centre and non-centre levels, and the BS.1770 weights are ignored.)

The same data bound what channel-wise attenuation can do for occupation. The passage has $817$ occupied seconds, $759$ of them in blocks. In the idealization, lowering the non-centre channels by $6$, $12$, or $20$ dB leaves $579$, $480$, or $441$ of the $759$ block seconds occupied, and removing them entirely leaves $435$ (about $57\%$). So even a perfect separation of the non-centre material would leave most of the blocks occupied: in $57.8\%$ of the passage's occupied seconds the non-centre channels are louder than the centre, yet the centre alone is still well above dialogue level. The loud material is not only in the surround and effects channels. This is why the monograph's downmix also compresses, and it shows that a stem-wise policy by channel is not by itself an occupation remedy for this passage.

\paragraph{7. Alternative edits.} Three transformations of increasing separation are on the monograph's own list: global gain, channel-based rebalancing with compression, and class-specific attenuation. Against the relations of Section~\ref{sec:altversions}:
\begin{center}
\small
\begin{tabular}{p{3.2cm}p{3.1cm}p{3.1cm}p{2.4cm}}
\toprule
 & occupation & $\rho_{\mathrm{dia}}$ & $\rho_{\mathrm{cue}}$\\
\midrule
global gain by $T_b$ & at least $95\%$ of block seconds cleared (computed) & fails by $T_b$ (computed) & not evaluable\\
channel rebalancing & idealized: removing all non-centre content leaves $435$ of $759$ block seconds occupied & median mix change $0.54$ dB in block speech, $\ge3.3$ dB in one in ten & not evaluable\\
class-specific & needs a detector with a stated error rate & depends on classifier & needs a cue list\\
\bottomrule
\end{tabular}
\end{center}
The compression stage of the channel rebalancing changes the mix everywhere and not only in blocks, as the monograph notes, so $\rho_{\mathrm{ord}}$-type relations on the whole film and not only on the passage would also have to be stated.

\paragraph{8. Evaluation.} The monograph specifies the study that would settle what the table cannot: logged volume behaviour under at least two playback conditions, one of them the dialogue-forward downmix, with comprehension questions after each viewing. In the vocabulary of this book, that study compares two versions that differ by a transformation whose contract is the list above, and it measures the preferred experience, the fourth claim of Section~\ref{sec:whisper}, which no computation here addresses.

\paragraph{What the real passage showed about the framework.}
\begin{itemize}
\item The occupation and recovery constructions transferred without change, and the monograph's definitions are refinements of them: a tolerance $\delta=10$ dB where the synthetic scene used $\mu=3$ dB, a merge gap, a minimum block length, and a recovery window of $5$ s at or below $D$. The propositions apply to the raw exceedance set, and statements about blocks must name the merge parameters.
\item The contract machinery was only half usable. $\rho_{\mathrm{dia}}$ could be assessed for global gain and, with the per-second table, for a channel policy, where it turned out that the dialogue inside blocks is far less centre-dominant than the film-wide figure suggests. $\rho_{\mathrm{cue}}$ and $\rho_{\mathrm{sal}}$ could not be assessed at all, because they depend on plot function, and the source records no plot.
\item The time-indexed analysis of Section~\ref{sec:timeindexed} had five columns. For this passage only observation and, partly, formal representation could be filled. Annotation, alternative edit, and evaluation needed information (shots, story events, listener data) that the source lacks. The synthetic scene looked easy to analyse in part because its events were stipulated together with their causal roles. Real passages do not come with that.
\item The response-model machinery of Chapter~\ref{ch:equivalence} had no data. One retrospective report is not a response model, so no sensitivity matrix or effective dimension can be estimated yet.
\item The monograph's grading of claims as established, measured, hypothesised, or rejected is the same discipline as Appendix~\ref{app:register}, and the transcribed numbers here are in its grade of measured.
\end{itemize}

\section{A shorter passage: speech timing against sound level}\label{sec:excerpt}

The previous section left the story and picture columns empty. This section takes a $130$ s excerpt inside the same passage, from $1{:}09{:}24$ to $1{:}11{:}34$ (seconds $4164$ to $4293$ of the per-second table), chosen because it lies wholly inside one occupation block and has a substantial number of subtitle-associated seconds. An earlier note gave this stretch by a clock reading that could be taken in two ways; the time here is hours, minutes and seconds from the start of the film. The section uses only the \emph{timing} of the subtitle cues, in the file \texttt{data/cues\_1h09\_1h11.csv}, and does not reproduce or interpret what is said. The subtitle-associated flag of the monograph marks a second when the cues cover at least half of it.

\paragraph{Observation.} All $130$ seconds lie in a block, and $106$ of them ($81.5\%$) are occupied under $\delta=10$ dB. The minimum exceedance over the excerpt is $5.3$ dB and the median $12.6$ dB, so no second of the excerpt is at or below the dialogue reference $D$, and no recovery window can occur. The flag marks $48$ seconds. They fall in short runs, with one long speech-free stretch: from $1{:}09{:}33$ the flag is absent for $41$ s, up to $1{:}10{:}13$. Table~\ref{tab:excerpt} gives ten-second segments.

\begin{table}[ht]
\centering\small
\caption{Excerpt $1{:}09{:}24$ to $1{:}11{:}34$ in ten-second segments, from the per-second table and the cue timings. ``Sub.'' is the number of subtitle-associated seconds, ``Occ.'' the number of occupied seconds, $E$ the median exceedance over $D=-34.6$ LUFS, and ``adv.'' the median centre advantage.}\label{tab:excerpt}
\begin{tabular}{lrrrr}
\toprule
segment start & Occ. & Sub. & median $E$ (dB) & median adv. (dB)\\
\midrule
$1{:}09{:}24$ & 8 & 3 & 10.9 & $-7.8$\\
$1{:}09{:}34$ & 10 & 0 & 15.1 & $-6.4$\\
$1{:}09{:}44$ & 9 & 0 & 14.7 & $-6.6$\\
$1{:}09{:}54$ & 7 & 0 & 12.4 & $-12.7$\\
$1{:}10{:}04$ & 10 & 0 & 13.2 & $-2.3$\\
$1{:}10{:}14$ & 10 & 10 & 12.9 & $10.4$\\
$1{:}10{:}24$ & 2 & 4 & 8.8 & $0.1$\\
$1{:}10{:}34$ & 8 & 5 & 13.9 & $9.0$\\
$1{:}10{:}44$ & 8 & 6 & 11.9 & $4.5$\\
$1{:}10{:}54$ & 9 & 6 & 11.8 & $6.2$\\
$1{:}11{:}04$ & 10 & 3 & 12.5 & $0.4$\\
$1{:}11{:}14$ & 8 & 8 & 11.2 & $8.2$\\
$1{:}11{:}24$ & 7 & 3 & 12.5 & $1.7$\\
\bottomrule
\end{tabular}
\end{table}

\paragraph{What the table shows.} Two features are visible without any interpretation of content. First, the exceedance does not follow speech density: the five segments with at most three flagged seconds have medians between $10.9$ and $15.1$ dB, and the speech-dense segments between $8.8$ and $13.9$ dB. Second, the centre advantage does: in the $41$ s speech-free stretch the segment medians are between $-12.7$ and $-6.4$ dB, so the non-centre channels dominate, and once speech begins the medians are mostly positive. Over the excerpt the median centre advantage is $+9.3$ dB in subtitle-associated seconds (three of the $48$ are negative) and $-4.8$ dB elsewhere. The acoustic organization of the excerpt is therefore a loud, non-centre-dominated stretch followed by a loud stretch in which speech is carried on top of material that remains well above $D$. The lowest exceedance (segment starting $1{:}10{:}24$, with only two occupied seconds) is a dip within the block, not a recovery, since it stays $5.3$ dB above $D$ at its minimum.

\paragraph{Which of the five columns can be filled.} Observation: yes, as above. Formal representation: partly. The excerpt is a path in the product $C\times A\times L\times E$ of Chapter~\ref{ch:scenario} whose acoustic coordinate is the exceedance sequence of Table~\ref{tab:excerpt}, and the occupation relation of Section~\ref{sec:occupation} holds on it for $106$ of $130$ seconds. Annotation, alternative edit and evaluation: no. They need what the source does not record: cuts, camera movement and visible action (for pacing); what is learned by whom (for narrative change); and which sounds are the only evidence of an event (for $\rho_{\mathrm{cue}}$). The next table lists, for this excerpt, exactly what a video excerpt would add. It is a worksheet, not a finding.

\begin{center}
\small
\begin{tabular}{p{3.6cm}p{4.2cm}p{4.2cm}}
\toprule
layer & question for the excerpt & source needed\\
\midrule
visual pacing & do cut rate and camera motion differ between the speech-free stretch and the speech-dense stretch? & shot boundaries, frame-differencing or manual log\\
narrative change & does the knowledge state (Chapter~\ref{ch:layers}) change at the start of the speech-dense stretch, or is it constant across both? & annotated story events with time stamps\\
acoustic organization & do occupation and exceedance track either of the above? & the table above (available)\\
\bottomrule
\end{tabular}
\end{center}

With these in hand the connection can be tested rather than asserted. If the shot-change rate is high in the speech-free stretch and the state layer is constant, the excerpt would exhibit visual and acoustic agitation without narrative change; if the state layer changes at the onset of speech while the exceedance does not, the acoustic organization is independent of the narrative change. Either outcome is a statement about one $130$ s excerpt, and neither bears on the film as a whole or on any viewer.


\section*{Exercises}

\begin{exercise}
Choose a film that uses concealed cuts and identify the seams. For each, estimate informally which of $\delta_\Sc$ and $\delta_\Img$ is small and state what kind of object conceals the discontinuity.
\end{exercise}

\begin{exercise}
For a film with an unreliable narrator, set up a sheaf-theoretic model with local sections given by what the narrator reports on each scene. State the condition under which the narrator's account is globally consistent and the condition under which it is not.
\end{exercise}

\begin{exercise}
Extend the synthetic scene by adding a third character $C$ who knows the letter is genuine. Write the story state, the typing of the reveals, and the number of valid orderings of three reveal events.
\end{exercise}

\begin{exercise}
Carry out the protocol of Section~\ref{sec:protocol} on a passage of a film of your choice, of at most two minutes, supplying edition and timecodes. Mark each claim as observation, reading, or hypothesis, and state one alternative edit with the relation it preserves and the relation it breaks.
\end{exercise}

\begin{exercise}
In the time-indexed analysis of Section~\ref{sec:timeindexed}, identify the segments whose annotation would change if the audience were told at $t=18$ that $A$ knows ($V_4$). Which formal objects in the table change, and which are unchanged?
\end{exercise}
